\scanpage{7}
vaut mieux dès lors exiger que ces \unsure{sommes}
soient \uline{homographiques} par morceaux\ldots
Dans le premier cas (exiger que les
disques ds le découpé soient ponctués\ldots)
on trouve, en supposant $\unsure{\Pi|\pi}_0(K_0)\longrightarrow\unsure{\Pi|\pi}_0(X_0)$
surjectif, \deleted{et \unsure{qu'il n'y ait} pas de bouts isolés\ldots}
nb de modules \inserted{réels}
\[
N=\left(\sum_{\substack{f\in\unsure{\Pi|\pi}_0(X\setminus K)=F\\f\text{ pas \uline{un} disque}}}
(6g(f)+4\nu(f)\unsure{-}6)\right)+\deleted{\unsure{\mu}}
\]
\deleted{$\mu$ étant le nb de ``faces'' $f$
qui sont des \deleted{disques}\inserted{couronnes}.}

\margin{On suppose $X_0$ connexe
$K$ connexe
$\lambda$ la connexité
de $K$}
\inserted{$\unsure{N}$ à $g\geq2$}
\[
3g-3\qquad\text{où}\quad
2-2g=\unsure{\chi}=\sum_{f\in F}(2-2g(f)-\nu(f))
\inserted{+1-\lambda}
\]
\[
\text{i.e.}\quad 2g-2=\sum_{f\in F}(2g(f)-2+\nu(f))
+(\lambda-1)
\]
\[
g\unsure{-|=}1=\frac12\sum_{f\in F}[2g(f)-2+\nu(f)-1]
+\frac{\unsure{3|1}}2(\lambda-1)
\]
\[
\begin{aligned}
3(g-1)&=\frac32\sum_{f\in F}[2g(f)-2+\nu(f)]
+\frac32(\lambda-1)\\
&=\sum_{f\in F}\left[3(g(f)-1)+\frac32\unsure{\nu}(f)\right]
+\frac32(\lambda-1)
\end{aligned}
\]
\[
(g\geq2)\qquad
N\unsure{-|=}(\unsure{2|3}g-3)
=\sum_{\substack{f\in F\\f\text{ pas un disque}}}
\left(3g(f)+\frac{\unsure{5}}2\nu(f)-3\right)
+\deleted{\sum_{\substack{f\in F\\f\text{ disque}}}\unsure{(}}
\frac32(\mu+\unsure{1-\lambda})\overset{?}{\geq}0
\]
$\mu$ étant le nb de $f$ qui
sont des disques

\[
(g=1)\qquad N-1=\sum_{\substack{f\text{ pas un}\\\text{disque}}}
\bigl(\underbrace{4\nu(f)-6}_{\unsure{\geq}2}\bigr)-1
\]
[NB $g(f)=0$ tjrs]
% archive foliation: 6
